%%%% Better Poster Skill style2 template.
%%%% Uses the same shared class but puts a hero visual and evidence cards in the central field.

\documentclass[a0paper,fleqn]{betterposter}

% Institution palette. Replace both colors with the same RGB triplet for
% single-color institutions, or use a coordinated secondary accent for
% dual-color institutions such as NUS blue plus orange.
\definecolor{institutionprimary}{RGB}{0,82,155}
\definecolor{institutionsecondary}{RGB}{239,124,0}
\institutionpalette{institutionprimary}{institutionsecondary}

\renewcommand{\maincolumnbackgroundcolor}{institutionprimary}
\renewcommand{\maincolumnfontcolor}{white}
\renewcommand{\fontsizemain}{\fontsize{74}{88}\selectfont}

\newcommand{\PosterClaimTitle}{A claim-first title states exactly what changed}
\newcommand{\PosterClaimLineOne}{The main result is explicit.}
\newcommand{\PosterClaimLineTwo}{The mechanism is visible.}
\newcommand{\PosterClaimLineThree}{The evidence is inspectable.}
\newcommand{\PosterAuthors}{First Author, Second Author, Senior Author}
\newcommand{\PosterAffiliations}{Lab or Department, University or Institute}
\newcommand{\InstitutionName}{University or Institute Name}
\newcommand{\HeroSubclaim}{One sentence tells the audience why the result matters and what to inspect first.}
\newcommand{\HeroFigure}{../figures/hero-teaser.pdf}
\newcommand{\PaperQRCode}{../figures/qr/01-paper.png}
\newcommand{\ScanIcon}{../assets/icons/scan-qr-code-white.png}
\newcommand{\InstitutionLogo}{../assets/institutions/current-logo.png}
\newcommand{\ContactLine}{author@example.com}

\begin{document}
\betterposter{
  % CENTER: claim + hero visual + evidence, with QR anchored near the bottom.
  \vfill
  {\RaggedRight\fontsizemain\bfseries
    {\color{white}\PosterClaimLineOne}\\[0.08em]
    {\color{white!68!black}\PosterClaimLineTwo}\\[0.08em]
    {\color{white!50!black}\PosterClaimLineThree}\par}
  \vspace{0.024\paperheight}
  {\fontsize{30}{38}\selectfont \HeroSubclaim\par}
  \vspace{0.022\paperheight}
  \heroimage{\HeroFigure}{The hero figure explains the mechanism and shows the strongest result.}
  \vspace{0.012\paperheight}
  \begin{center}
    \evidencecard{What to see first}{Point to the visual pattern, comparison, or mechanism.}
    \hfill
    \evidencecard{Why it is credible}{Name the strongest validation, uncertainty check, or ablation.}
  \end{center}
  \vfill
  \inlineqrcodewithicon{\PaperQRCode}{\ScanIcon}{\textbf{Scan for paper, code, and supplemental figures.}}
}{
  % LEFT: context before the claim.
  \title{\PosterClaimTitle}
  \author{\PosterAuthors}
  \institution{\PosterAffiliations}

  \section{Why it matters}
  \begin{itemize}
    \item The real problem in one line.
    \item What current methods miss or cost.
    \item Who should care about the result.
  \end{itemize}

  \section{Method}
  \begin{enumerate}
    \item Input or experimental setup.
    \item Core transformation or intervention.
    \item Output, measurement, or decision.
  \end{enumerate}

  \section{Design choices}
  Keep only details needed to trust the central claim.

  \vfill
  \institutionbrand{\InstitutionLogo}{\InstitutionName}{0.62\textwidth}
  \posterfootnote{Contact: \ContactLine}
}{
  % RIGHT: technical conversation layer.
  \section{Top evidence}
  \begin{itemize}
    \item Evidence bullet with dataset, metric, and baseline.
    \item Qualitative example or ablation result.
    \item Robustness, uncertainty, or transfer test.
  \end{itemize}

  \section{Limitations}
  \begin{itemize}
    \item State the most important boundary honestly.
    \item Name where transfer is unknown.
    \item Point to the next critical experiment.
  \end{itemize}

  \section{Extra figure}
  \begin{center}
    \postergraphic[width=\textwidth]{../figures/extra-result.pdf}{Replace with a readable extra figure.}
  \end{center}

  \section{Scan}
  The QR code links to details that should not be forced into the poster.
}
\end{document}
